\documentclass[10pt,a4paper]{amsart}
\usepackage{amsmath,amssymb,amsthm,graphicx}
\usepackage{verbatim}
\newenvironment{reference}{\comment}{\endcomment}
%\newenvironment{reference}{}{}
\newenvironment{slogan}{\comment}{\endcomment}
\newenvironment{history}{\comment}{\endcomment}

% For commutative diagrams we use Xy-pic
\usepackage[all]{xy}

% We use 2cell for 2-commutative diagrams.
\xyoption{2cell}
% Theorem environments.
%
\theoremstyle{plain}
\newtheorem{theorem}[subsection]{Theorem}
\newtheorem{proposition}[subsection]{Proposition}
\newtheorem{lemma}[subsection]{Lemma}

\theoremstyle{definition}
\newtheorem{definition}[subsection]{Definition}
\newtheorem{example}[subsection]{Example}
\newtheorem{exercise}[subsection]{Exercise}
\newtheorem{situation}[subsection]{Situation}

\theoremstyle{remark}
\newtheorem{remark}[subsection]{Remark}
\newtheorem{remarks}[subsection]{Remarks}

\numberwithin{equation}{subsection}

% Macros
%
\def\lim{\mathop{\mathrm{lim}}\nolimits}
\def\colim{\mathop{\mathrm{colim}}\nolimits}
\def\Spec{\mathop{\mathrm{Spec}}}
\def\Hom{\mathop{\mathrm{Hom}}\nolimits}
\def\Ext{\mathop{\mathrm{Ext}}\nolimits}
\def\SheafHom{\mathop{\mathcal{H}\!\mathit{om}}\nolimits}
\def\SheafExt{\mathop{\mathcal{E}\!\mathit{xt}}\nolimits}
\def\Sch{\mathit{Sch}}
\def\Mor{\mathop{\mathrm{Mor}}\nolimits}
\def\Ob{\mathop{\mathrm{Ob}}\nolimits}
\def\Sh{\mathop{\mathit{Sh}}\nolimits}
\def\NL{\mathop{N\!L}\nolimits}
\def\CH{\mathop{\mathrm{CH}}\nolimits}
\def\proetale{{pro\text{-}\acute{e}tale}}
\def\etale{{\acute{e}tale}}
\def\QCoh{\mathit{QCoh}}
\def\Ker{\mathop{\mathrm{Ker}}}
\def\Im{\mathop{\mathrm{Im}}}
\def\Coker{\mathop{\mathrm{Coker}}}
\def\Coim{\mathop{\mathrm{Coim}}}

% Boxtimes
%
\DeclareMathSymbol{\boxtimes}{\mathbin}{AMSa}{"02}

%
% Macros for moduli stacks/spaces
%
\def\QCohstack{\mathcal{QC}\!\mathit{oh}}
\def\Cohstack{\mathcal{C}\!\mathit{oh}}
\def\Spacesstack{\mathcal{S}\!\mathit{paces}}
\def\Quotfunctor{\mathrm{Quot}}
\def\Hilbfunctor{\mathrm{Hilb}}
\def\Curvesstack{\mathcal{C}\!\mathit{urves}}
\def\Polarizedstack{\mathcal{P}\!\mathit{olarized}}
\def\Complexesstack{\mathcal{C}\!\mathit{omplexes}}
% \Pic is the operator that assigns to X its picard group, usage \Pic(X)
% \Picardstack_{X/B} denotes the Picard stack of X over B
% \Picardfunctor_{X/B} denotes the Picard functor of X over B
\def\Pic{\mathop{\mathrm{Pic}}\nolimits}
\def\Picardstack{\mathcal{P}\!\mathit{ic}}
\def\Picardfunctor{\mathrm{Pic}}
\def\Deformationcategory{\mathcal{D}\!\mathit{ef}}
\usepackage{cleveref}
\crefname{originalsection}{original section}{original sections}
\Crefname{originalsection}{Original section}{Original sections}
\makeatletter
\newlabel{section-localize}{{03DH}{source}{}{sourcepointer-1}{}}
\newlabel{section-localize@cref}{{[originalsection][0][]03DH}{[1][0][]source}{}{}{}}
\makeatother
\crefname{originalreference}{original reference}{original references}
\Crefname{originalreference}{Original reference}{Original references}
\makeatletter
\newlabel{sites-lemma-localize-topos}{{04GZ}{source}{}{sourcepointer-2}{}}
\newlabel{sites-lemma-localize-topos@cref}{{[originalreference][0][]04GZ}{[1][0][]source}{}{}{}}
\makeatother
\crefname{originallemma}{original lemma}{original lemmas}
\Crefname{originallemma}{Original lemma}{Original lemmas}
\makeatletter
\newlabel{lemma-extension-by-zero}{{03DI}{source}{}{sourcepointer-3}{}}
\newlabel{lemma-extension-by-zero@cref}{{[originallemma][0][]03DI}{[1][0][]source}{}{}{}}
\makeatother
\title{The localization of a ringed topos at a sheaf is a ringed topos with a canonical morphism to the original}
\author{The Stacks Project Authors}
\date{stacks-project:sites-modules.tex:local-snapshot-20261012; GNU FDL 1.2 or later}
\begin{document}
\maketitle
\noindent\textit{Editorial note.} All delivered mathematical text is the author's own, copied verbatim from the pinned Stacks Project chapter file sites-modules.tex; no statement or proof step was rewritten, reordered or removed.\par\medskip
\noindent\textit{Editorial note.} The counted span cites one sites-level lemma, recorded in assembly.external\_references with its canonical tag as a navigation pointer.\par\medskip
\noindent\textit{Allowed standard prerequisites (editorial).}\begin{itemize}
\item Localization of a topos at a sheaf. Stacks Project, Sites, Lemma sites-lemma-localize-topos (tag 00YV). Supplies the topos structure of the localization and the reduction to the representable case.
\end{itemize}\medskip
\ifdefined\hypertarget\hypertarget{sourcepointer-1}{}\fi
\noindent\textit{Source reference.} Original section 03DH: sources/stacks-chow/tags.txt, source lines 4420--4420. Cross-chapter standard prerequisite cited by the author inside the counted span; it is used as a pointer and its own proof is not part of this item.\par\medskip
\ifdefined\hypertarget\hypertarget{sourcepointer-2}{}\fi
\noindent\textit{Source reference.} Original reference 04GZ: sources/stacks-chow/tags.txt, source lines 5850--5850. Cross-chapter standard prerequisite cited by the author inside the counted span; it is used as a pointer and its own proof is not part of this item.\par\medskip
\ifdefined\hypertarget\hypertarget{sourcepointer-3}{}\fi
\noindent\textit{Source reference.} Original lemma 03DI: sources/stacks-chow/tags.txt, source lines 4421--4421. Cross-chapter standard prerequisite cited by the author inside the counted span; it is used as a pointer and its own proof is not part of this item.\par\medskip
\begin{lemma}
\label{lemma-localize-ringed-topos}
Let $(\Sh(\mathcal{C}), \mathcal{O})$ be a ringed topos.
Let $\mathcal{F} \in \Sh(\mathcal{C})$ be a sheaf.
For a sheaf $\mathcal{H}$ on $\mathcal{C}$ denote
$\mathcal{H}_\mathcal{F}$ the sheaf $\mathcal{H} \times \mathcal{F}$
seen as an object of the category $\Sh(\mathcal{C})/\mathcal{F}$.
The pair
$(\Sh(\mathcal{C})/\mathcal{F}, \mathcal{O}_\mathcal{F})$
is a ringed topos and there is a canonical morphism of ringed topoi
$$
(j_\mathcal{F}, j_\mathcal{F}^\sharp) :
(\Sh(\mathcal{C})/\mathcal{F}, \mathcal{O}_\mathcal{F})
\longrightarrow
(\Sh(\mathcal{C}), \mathcal{O})
$$
which is a localization as in
Section \ref{section-localize}
such that
\begin{enumerate}
\item the functor $j_\mathcal{F}^{-1}$ is the functor
$\mathcal{H} \mapsto \mathcal{H}_\mathcal{F}$,
\item the functor $j_\mathcal{F}^*$ is the functor
$\mathcal{H} \mapsto \mathcal{H}_\mathcal{F}$,
\item the functor $j_{\mathcal{F}!}$ on sheaves of sets is the forgetful
functor $\mathcal{G}/\mathcal{F} \mapsto \mathcal{G}$,
\item the functor $j_{\mathcal{F}!}$ on sheaves of modules associates
to the $\mathcal{O}_\mathcal{F}$-module
$\varphi : \mathcal{G} \to \mathcal{F}$ the $\mathcal{O}$-module
which is the sheafification of the presheaf
$$
V \longmapsto
\bigoplus\nolimits_{s \in \mathcal{F}(V)}
\{\sigma \in \mathcal{G}(V) \mid \varphi(\sigma) = s \}
$$
for $V \in \Ob(\mathcal{C})$.
\end{enumerate}
\end{lemma}

\begin{proof}
By
Sites, Lemma \ref{sites-lemma-localize-topos}
we see that $\Sh(\mathcal{C})/\mathcal{F}$ is a topos
and that (1) and (3) are true. In particular this shows that
$j_\mathcal{F}^{-1}\mathcal{O} = \mathcal{O}_\mathcal{F}$
and shows that $\mathcal{O}_\mathcal{F}$ is a sheaf of rings.
Thus we may choose the map $j_\mathcal{F}^\sharp$ to be the identity,
in particular we see that (2) is true.
Moreover, the proof of
Sites, Lemma \ref{sites-lemma-localize-topos}
shows that we may assume $\mathcal{C}$ is a site with all finite limits
and a subcanonical topology and that $\mathcal{F} = h_U$ for some object
$U$ of $\mathcal{C}$.
Then (4) follows from the description of $j_{\mathcal{F}!}$ in
Lemma \ref{lemma-extension-by-zero}.
Alternatively one could show directly that the functor described
in (4) is a left adjoint to $j_\mathcal{F}^*$.
\end{proof}

\end{document}
